\paragraph{\trNaturalTransformation。}\label{sub:natural-transformations}
\trNonunital{} \trAFunctor{} $\scr{F}\colon\scr{A}\to\scr{B}$ 自身就构成某个\trNonunital{}
\trACategory{} $\scr{Q}=\mathit{nu}\text{-}\mathit{fun}(\scr{A},\scr{B})$
中的\trObject{}，而这个\trCategory{}中\trMorphism{}链空间的一个元素
$T\in\hom^{g}_{\scr{Q}}(\scr{F}_0,\scr{F}_1)$\footnote{译注：原书此处印作斜体 $hom$；
疑应为正体 $\hom$（原书同位置记号恒排正体，十二处抽验仅此一处排斜体，系孤例）；
此处径改。——译者注} 是一个序列 $(T^0,T^1,\ldots)$，
其中每个 $T^d$ 都是一族\trMultilinear{}映射
\begin{equation*}
\hom_{\scr{A}}(X_{d-1},X_d)\otimes\cdots\otimes\hom_{\scr{A}}(X_0,X_1)
  \longrightarrow\hom_{\scr{B}}(\scr{F}_0X_0,\scr{F}_1X_d)[g-d],
\end{equation*}
对一切 $(X_0,\ldots,X_d)$。特别地，$T^0$ 是
$\hom^{g}_{\scr{B}}(\scr{F}_0X,\scr{F}_1X)$ 中的一族元素，
$\scr{A}$ 的每个\trObject{} $X$ 对应其一。遵循 \trFukaya，我们把这样的 $T$ 称为从
$\scr{F}_0$ 到 $\scr{F}_1$ 的\textbf{预\trNaturalTransformation}（pre-natural
transformation），其\trDegree{}为 $g=|T|$。\textbf{\trBoundaryOperator}（boundary
operator）为
\begin{equation}\label{eq:natural-transformation-boundary}
\begin{split}
\mu^1_{\scr{Q}}(T)^d(a_d,\ldots,a_1)
&=\sum_{r,i}\sum_{s_1,\ldots,s_r}(-1)^{\dagger}\mu^r_{\scr{B}}\bigl(
  \scr{F}_1^{s_r}(a_d,\ldots,a_{d-s_r+1}),\ldots,\\
&\qquad \scr{F}_1^{s_i+1}(\ldots,a_{s_1+\cdots+s_i+1}),\\
&\qquad T^{s_i}(a_{s_1+\cdots+s_i},\ldots,a_{s_1+\cdots+s_{i-1}+1}),\\
&\qquad \scr{F}_0^{s_{i-1}}(a_{s_1+\cdots+s_{i-1}},\ldots),\ldots,
  \scr{F}_0^{s_1}(a_{s_1},\ldots,a_1)\bigr)\\
&\quad-\sum_{m,n}(-1)^{\clubsuit_n+|T|-1}\,T^{d-m+1}(a_d,\ldots,a_{n+m+1},\\
&\qquad \mu^m_{\scr{A}}(a_{n+m},\ldots,a_{n+1}),a_n\ldots,a_1).
\end{split}
\end{equation}
第一个求和取遍 $1\le i\le r$ 与\trPartition{} $s_1+\cdots+s_r=d$，其中诸 $s_i$ 容许为零；而
$\dagger=(|T|-1)(|a_1|+\cdots+|a_{s_1+\cdots+s_{i-1}}|-s_1-\cdots-s_{i-1})$。
$\mu^1_{\scr{Q}}(T)=0$ 的\trCocycle{}恰是真正的\textbf{\trNaturalTransformation}（natural
transformation），而\trCoboundary{}则给出它们之间\trChainHomotopy{}的概念。至于
$T_1\in\hom_{\scr{Q}}(\scr{F}_0,\scr{F}_1)$ 与
$T_2\in\hom_{\scr{Q}}(\scr{F}_1,\scr{F}_2)$ 的\trComposition{}公式则要简单得多：
\begin{equation}\label{eq:q-composition-formula}
\begin{split}
\mu^2_{\scr{Q}}(T_2,T_1)^d(a_d,\ldots,a_1)
&=\sum_{r,i,j}\sum_{s_1,\ldots,s_r}(-1)^{\circ}\,\mu^r_{\scr{B}}\bigl(
  \scr{F}_2^{s_r}(a_d,\ldots,a_{d-s_r+1}),\ldots,\\
&\qquad \scr{F}_2^{s_j+1}(\ldots),\\
&\qquad T_2^{s_j}(a_{s_1+\cdots+s_j},\ldots,a_{s_1+\cdots+s_{j-1}+1}),\\
&\qquad \scr{F}_1^{s_{j-1}}(\ldots),\ldots,\scr{F}_1^{s_i+1}(\ldots),\\
&\qquad T_1^{s_i}(a_{s_1+\cdots+s_i},\ldots,a_{s_1+\cdots+s_{i-1}+1}),\\
&\qquad \scr{F}_0^{s_{i-1}}(\ldots),\ldots,\scr{F}_0^{s_1}(a_{s_1},\ldots,a_1)\bigr),
\end{split}
\end{equation}
其中
$\circ=\sum\limits_{k=1}^{s_1+\cdots+s_{j-1}}(|T_2|-1)(|a_k|-1)
+\sum\limits_{k=1}^{s_1+\cdots+s_{i-1}}(|T_1|-1)(|a_k|-1)$。注意 $\mu^{\bullet}_{\scr{A}}$
不再出现；$\mu^d_{\scr{Q}}$（$d>2$）的公式完全遵循\eqref{eq:q-composition-formula}
的同一模式。

设 $T$ 是介于 \trAFunctor{} $\scr{F}_0,\scr{F}_1\colon\scr{A}\to\scr{B}$ 之间的一个\trNaturalTransformation{}，考察元素
$[T^0_X]\in\Hom_{H(\scr{B})}(\scr{F}_0X,\scr{F}_1X)$。它们满足如下的自然性条件：
\begin{equation*}
[T^0_{X^1}]\cdot[\scr{F}_0^1(a)]=(-1)^{|a|\cdot|T|}[\scr{F}_1^1(a)]\cdot[T^0_{X^0}]
\end{equation*}
其中 $[a]\in\Hom_{H(\scr{A})}(X^0,X^1)$。于是，它们构成介于\trCohomology\trFunctor{}%
$H(\scr{F}_0)$ 与 $H(\scr{F}_1)$ 之间的一个\trNaturalTransformation{}，记作 $H(T)$。
更正式地，把由\trNonunital（\trLinear、\trGraded）\trFunctor{}%
$H(\scr{A})\to H(\scr{B})$ 及其间的\trNaturalTransformation{}所组成的\trCategory{}%
记为 $\mathit{Nu}\text{-}\mathit{Fun}(H(\scr{A}),H(\scr{B}))$。那么，指定
$\scr{F}\mapsto H(\scr{F})$ 与 $T\mapsto H(T)$ 便定义了一个\trNonunital\trFunctor
\begin{equation}\label{eq:H-of-nu-fun-to-Nu-Fun}
H(\mathit{nu}\text{-}\mathit{fun}(\scr{A},\scr{B}))
  \longrightarrow\mathit{Nu}\text{-}\mathit{Fun}(H(\scr{A}),H(\scr{B})).
\end{equation}
